% Classe de document et géométrie
\documentclass[a4paper,10pt]{article}
\usepackage[utf8]{inputenc}
\usepackage[T1]{fontenc}
\usepackage{geometry}
\geometry{left=1cm,right=1cm,top=1.2cm,bottom=1.2cm}

% Paquets
\usepackage{parskip}
\usepackage{enumitem}
\usepackage{hyperref}
\usepackage{titlesec}
\usepackage{xcolor}
\usepackage{FiraSans}
\renewcommand{\familydefault}{\sfdefault}

% Configuration des liens
\hypersetup{
    colorlinks=true,
    linkcolor=blue,
    urlcolor=blue
}

% Titres de section
\titleformat{\section}{\large\bfseries}{}{0em}{}[\titlerule]
\titlespacing{\section}{0pt}{0.6em}{0.4em}

\raggedbottom

% Commande pour l'en-tête d'expérience
\newcommand{\jobheader}[4]{%
    \textbf{#1, #2} \\
    \textit{#3, #4}
    \vspace{0.05cm}
}

\begin{document}

% En-tête
    {\LARGE\bfseries Rachid Rodrigue BADINI} \\[0.15cm]
    \textbf{Ingénieur Fullstack Senior | React \textbullet\, Next.js \textbullet\, TypeScript \textbullet\, Python \textbullet\, Node.js} \\
    {\footnotesize Téléphone : +226 75 22 53 88 | Courriel : bsrodrigue@gmail.com | LinkedIn : \href{https://linkedin.com/in/b-rodrigue}{linkedin.com/in/b-rodrigue} | GitHub : \href{https://github.com/bsrodrigue}{github.com/bsrodrigue} \\
    \textbf{Lieu :} Burkina Faso, Ouagadougou}

\vspace{0.1cm}

% Résumé professionnel
\section*{Résumé Professionnel}
Ingénieur Fullstack avec plus de 5 ans d'expérience dans la conception et le développement d'applications web scalables, de systèmes backend distribués et d'interfaces utilisateur haute performance avec Python, Django, FastAPI, React, Next.js, TypeScript, PostgreSQL, Redis, Docker et Kubernetes. Expérimenté dans l'amélioration des performances applicatives, l'optimisation de bases de données et la conception d'API fiables pour des environnements de production. Combine une connaissance approfondie de l'infrastructure backend avec le développement frontend moderne pour fournir des solutions de bout en bout. L'objectif à long terme est de construire des plateformes haute performance et des infrastructures pour développeurs utilisées à l'échelle mondiale.

% Compétences techniques
\section*{Compétences Techniques}
\begin{itemize}[leftmargin=*,nosep]
    \item \textbf{Langages :} Python (Expert), Go (Avancé), C++ (Systèmes/Qt), SQL (PostgreSQL/MySQL), TypeScript, JavaScript (ES6+), Dart
    \item \textbf{Frontend :} React (v17+), Next.js, TypeScript, Apollo GraphQL, Hooks, Material UI, Formik, Tailwind CSS, React Native, Flutter
    \item \textbf{Backend et Systèmes :} Systèmes distribués, Microservices, Node.js, NestJS, Architecture Modulaire, Concurrency, Optimisation de performance, Débogage et Profilage
    \item \textbf{Infrastructure et DevOps :} AWS (EC2, S3, Lambda, ECS, Fargate, RDS), Docker, Kubernetes, Ansible, Terraform, Prometheus, Grafana, CI/CD (GitHub Actions)
    \item \textbf{Bases de données et Cache :} PostgreSQL (Optimisation de requêtes, Indexation), MySQL, Redis, Files d'attente (Celery/RabbitMQ)
    \item \textbf{Mobile :} React Native, Flutter, Gestion d'État (Redux, BLoC, Riverpod), Firebase, Fastlane
    \item \textbf{Tests et Outils :} Jest, Cypress, Detox, PyTest, GitFlow, Jira, Scrum
    \item \textbf{Outils IA :} Claude Code, OpenCode, Gemini CLI
\end{itemize}

% Expérience professionnelle
\section*{Expérience Professionnelle}

\jobheader{Ingénieur Fullstack Principal}{DISCOM}{Mars 2026 -- Présent}{Ouagadougou}
\begin{itemize}[leftmargin=*,nosep]
    \item \textbf{Architecture Multi-Locataire (Multi-Tenant) :} Conception d'une application critique pour une structure nationale majeure (AICB - Association Interprofessionnelle du Coton du Burkina), en respectant rigoureusement les meilleures pratiques Django et le Clean Code.
    \item \textbf{Formulaires Dynamiques et Ingestion Massive :} Développement d'un système robuste de gestion de formulaires dynamiques supportant des imports CSV massifs et concurrents (traitement de plus de 200 000 lignes par fichier).
    \item \textbf{Systèmes à haut débit :} Identification des goulots d'étranglement PostgreSQL (pooling de connexions et requêtes N+1) dans le pipeline d'ingestion CSV. Refactorisation des requêtes et optimisation du pool de connexions, réduisant l'ingestion de 10 000 lignes de 67s à 0,7s (amélioration de 96\%) et traitant 185 862 lignes en 14s sans accumulation.
    \item \textbf{Architecture Asynchrone :} Mise en œuvre de Redis et Celery pour le traitement des tâches en arrière-plan, découplant avec succès les opérations de données longues du cycle requête-réponse.
    \item \textbf{Observabilité et Monitoring :} Développement d'un code d'instrumentation personnalisé pour l'ingestion CSV exposé via un endpoint /metrics et création de tableaux de bord Grafana sur mesure pour obtenir des pistes d'optimisation exploitables.
    \item \textbf{Analyse des Goulots d'Étranglement BD :} Utilisation de pg\_stat\_statements avec psql pour identifier et résoudre les goulots d'Étranglement PostgreSQL dans la pipeline de traitement.
    \item \textbf{Communication avec les Parties Prenantes :} Principal point de contact avec les parties prenantes du projet à l'AICB, assurant la collecte des besoins et le conseil technique.
\end{itemize}

\vspace{0.15cm}

\jobheader{Ingénieur Logiciel Fullstack (C++)}{Astek Madagascar}{Juillet 2025 -- Décembre 2025}{Madagascar}
\begin{itemize}[leftmargin=*,nosep]
    \item \textbf{Logiciel Système Bas Niveau :} Développement de logiciels C/C++ robustes, multithreadés et modulaires pour des équipements de billetterie fonctionnant dans des environnements contraints.
    \item \textbf{Implémentation de Fonctionnalités et UI :} Implémentation de nouvelles fonctionnalités et interfaces utilisateur conformément à des documents de spécification stricts.
    \item \textbf{Débogage de Protocole Propriétaire :} Apprentissage et débogage de requêtes réseau basées sur un protocole réseau propriétaire.
    \item \textbf{Collaboration à la Conception :} Discussion et affinement des choix de conception avec les membres de l'équipe pour garantir des solutions robustes et maintenables.
    \item \textbf{Correction de Bogues :} Résolution de bogues d'interface et de bas niveau pour améliorer la stabilité du système et l'expérience utilisateur.
    \item \textbf{Cloud et Outils :} Utilisation de Microsoft Azure pour les déploiements et le suivi des bogues.
\end{itemize}

\vspace{0.15cm}

\jobheader{Ingénieur Fullstack Principal (Next.js \& Python)}{Doonya Labs}{Novembre 2024 -- Juillet 2025}{Hybride}
\begin{itemize}[leftmargin=*,nosep]
    \item \textbf{Scalabilité :} Conception d'une plateforme Django/DRF mise à l'échelle à plus de 50 000 requêtes API quotidiennes avec des temps de réponse P95 inférieurs à 200ms.
    \item \textbf{Développement Frontend :} Création d'interfaces réactives avec React/Next.js et gestion des transitions d'état complexes via les Hooks et les bibliothèques de formulaires modernes.
    \item \textbf{Infrastructure :} Provisionnement de l'infrastructure AWS avec Terraform et orchestration de conteneurs via Kubernetes. Mise en place d'un pipeline de déploiement GitHub Actions pour des mises en production automatisées et reproductibles.
    \item \textbf{Observabilité :} Mise en place de Prometheus et Grafana pour le monitoring et les métriques, assurant une visibilité en temps réel sur les performances du système.
    \item \textbf{Leadership :} Mentorat d'une équipe pluridisciplinaire de 4 ingénieurs et direction des cérémonies SCRUM.
\end{itemize}

\vspace{0.15cm}

\jobheader{Ingénieur Fullstack Principal (React \& React Native \& Python)}{GandyamLigdi}{Juillet 2023 -- Octobre 2024}{À distance}
\begin{itemize}[leftmargin=*,nosep]
    \item \textbf{Systèmes financiers :} Développement d'API financières sécurisées et conformes ACID traitant plus de 10 000 transactions mensuelles.
    \item \textbf{Livraison Fullstack :} Développement de fonctionnalités sur le web (React) et mobile (React Native/TypeScript), avec un focus sur l'optimisation des performances et l'itération rapide.
    \item \textbf{Architecture :} Conception d'un monolithe modulaire découplé utilisant le "Repository Pattern" et l'injection de dépendances.
    \item \textbf{Optimisation :} Augmentation du débit de 40 \% via une stratégie de mise en cache Redis et l'optimisation des requêtes PostgreSQL.
    \item \textbf{CI/CD :} Configuration d'un pipeline CI/CD avec GitHub Actions pour automatiser les déploiements et garantir la reproductibilité des mises en production.
\end{itemize}

\vspace{0.15cm}

\jobheader{Consultant Fullstack}{Freelance}{2019 -- Présent}{À distance}
\begin{itemize}[leftmargin=*,nosep]
    \item \textbf{PMUB (React Native) :} Développement d'une plateforme de paris hippiques gérant des flux de données en temps réel via WebSockets et intégrations de paiement sécurisées.
    \item \textbf{Songre - EliteApp (React Native) :} Architecture d'une super-app multi-services (Emploi, E-commerce, Médical, Livraison) gérant plusieurs types d'utilisateurs (clients, vendeurs, livreurs, admins).
    \item \textbf{Allo Youri (Flutter \& Web) :} Implémentation d'une plateforme de livraison avec suivi GPS en temps réel (Maps API) et création d'une landing page responsive.
    \item \textbf{Les Genets (Flutter) :} Développement d'une application médicale offrant téléconsultations, services SOS et gestion de dossiers médicaux.
    \item \textbf{AppShare (React Native \& Go) :} Conception d'un outil de distribution d'APK avec un backend robuste en Go, optimisant le versioning et le partage de binaires.
    \item \textbf{Dom Immo (NextJS) :} Développement de l'interface utilisateur haute performance pour une plateforme immobilière propulsée par l'IA.
\end{itemize}

\vspace{0.15cm}

\jobheader{Ingénieur Logiciel Fullstack (React Native • Node • Python)}{N7 Studio}{Avril 2021 -- Avril 2023}{À distance (Canada)}
\begin{itemize}[leftmargin=*,nosep]
    \item \textbf{Haute concurrence :} Développement de services Node.js + TypeScript + GraphQL pour une plateforme de streaming supportant 5 000 utilisateurs simultanés.
    \item \textbf{Développement Frontend :} Développement d'interfaces web réactives avec React et d'applications mobiles utilisant React Native et TypeScript.
    \item \textbf{Fiabilité :} Intégration de suites de tests complètes (PyTest, Cypress, Jest) dans les pipelines automatisés, atteignant 85 \% de couverture de code.
    \item \textbf{Efficacité :} Réduction des temps de réponse API de 45 \% par l'ajustement de la taille du pool de connexions et l'optimisation des requêtes.
\end{itemize}

\vspace{0.15cm}

\jobheader{Administrateur Systèmes Linux (Stage)}{ANPTIC}{Mars 2021 -- Juin 2021}{Ouagadougou}
\begin{itemize}[leftmargin=*,nosep]
    \item \textbf{Optimisation Noyau :} Personnalisation et recompilation de noyaux Linux pour supporter des pilotes spécifiques et renforcer la sécurité.
    \item \textbf{Monitoring Réseau :} Surveillance d'un réseau national avec NAGIOS pour garantir la disponibilité de l'infrastructure et une réponse rapide aux incidents.
\end{itemize}

% Formation
\section*{Formation}
\textbf{Maîtrise (M.Sc.) en Génie Logiciel} \hfill Université Aube Nouvelle \\
\textit{Spécialisation en Systèmes Distribués et Architecture} \hfill 2021 -- 2022

\vspace{0.1cm}

\textbf{Licence (B.Sc.) en Réseaux Informatiques} \hfill Université Aube Nouvelle \\
\textit{Focus sur TCP/IP, Systèmes Linux et Administration Réseau} \hfill 2017 -- 2021

% Projets d'Ingénierie Clés
\section*{Projets}
\begin{itemize}[leftmargin=*,nosep]
    \item \textbf{DreamServer (Go) :} Conception d'un reverse proxy HTTP et d'un serveur de fichiers statiques à partir de zéro. Développement d'un parseur HTTP personnalisé pour gérer les connexions TCP brutes, le transfert d'en-têtes et les journaux structurés avec suivi de latence. \\
    GitHub : \href{https://github.com/bsrodrigue/dreamproxy}{github.com/bsrodrigue/dreamproxy} (14 stars)
    \item \textbf{Buildshare :} Un système full-stack Django/Celery permettant aux équipes de développement d'importer des binaires APK, de les organiser par projet et de distribuer des versions publiées aux testeurs via un client mobile dédié. \\
    GitHub : \href{https://github.com/bsrodrigue/buildshare}{github.com/bsrodrigue/buildshare} (26 stars)
    \item \textbf{Memscan (C) :} Un petit scanner mémoire en C pour inspecter et manipuler la mémoire des processus Linux. \\
    GitHub : \href{https://github.com/bsrodrigue/memscan}{github.com/bsrodrigue/memscan}
\end{itemize}

\vspace{0.1cm}

% Langues
\section*{Langues}
\textbf{Français :} Maternel | \textbf{Anglais :} C1 (Courant professionnel)

\end{document}
